\section{Introduction}
\label{sec:intro}

Understanding 3D geometry is a fundamental challenge in computer vision, central to fields ranging from robotics and AR/VR to content creation and scientific analysis.

Recent efforts to build 3D foundation models~\cite{xue2023ulip, liu2023openshape, zhou2023uni3d} have achieved impressive results on \emph{global} tasks like retrieval and classification, surpassing by a large margin VLM-based approaches such as PointCLIPv2 \cite{zhu2022pointclip}, which rely solely on 2D renderings for 3D understanding. However, open-vocabulary dense prediction tasks, such as 3D part segmentation, remain largely dominated by multi-view pipelines~\cite{zhu2022pointclip, liu2023partslip, abdelreheem2023satr, garosi20253d}. These methods render each 3D shape into multiple images, extract features using powerful 2D vision models like DINOv2~\cite{oquab2023dinov2} and CLIP~\cite{radford2021learning}, and then fuse the resulting 2D predictions back into 3D. This paradigm effectively transfers informative 2D semantic priors to 3D, enabling open-vocabulary 3D part understanding, and has therefore become the \emph{de facto} approach for dense 3D tasks.

Despite being effective, multi-view pipelines lack geometric grounding, as their predictions are primarily based on 2D appearance cues rather than the underlying 3D structure. Their inference process is also computationally expensive,  

due to the need for extensive multi-view rendering and per-view inference, followed by complex geometric fusion. 

Finally, their performance relies heavily on LLM-driven prompt engineering on test sets, and they degrade significantly when presented with generic part labels common in real-world applications~\cite{garosi20253d}.



We propose \textbf{PatchAlign3D}, the first encoder-only 3D model that learns language-aligned \emph{local} features directly from point clouds. 
Our model bypasses the limitations of multi-view pipelines, achieving high-performance open-world 3D part segmentation in a single feed-forward pass.
We introduce a two-stage training strategy. Stage 1 performs 2D-to-3D feature distillation using a 3D transformer encoder, transferring dense visual features from a pre-trained 2D model such as DINOv2 \cite{oquab2023dinov2} to 3D patch tokens. This step lays the foundation for equipping 3D geometric encoding with fine-grained visual representations. Stage 2 then aligns these 3D patch embeddings with textual part descriptions encoded by a pre-trained text encoder, such as CLIP \cite{radford2021learning}.

A central challenge in our approach is learning from large-scale 3D shape segmentation data. We leverage 3D part annotations from Find3D \cite{ma2025find}, a recent data engine that automatically segments and annotates shape parts through a 2D segmentor SAM\cite{kirillov2023segment} and a VLM~\cite{team2023gemini}. These annotations are, however, inherently noisy and inconsistent across parts, for example, a single 3D patch may be associated with multiple part names, and segmentation masks are often fragmented or incomplete.
Key to our approach is that, instead of learning on unreliable point-level annotations, we perform semantic alignment at the \emph{patch level}. This local aggregation averages out annotation noise and is more robust to inconsistent boundaries. Furthermore, we propose a multi-positive sample-wise contrastive objective that uses fractional labels to handle ambiguous segmentation, yielding robust, generalizable geometric representations.



In summary, our contributions are: 
\begin{itemize}
    \item We introduce the first 3D encoder that produces language-aligned, patch-level features, closing the gap between global 3D foundation models and multi-view dense VLM pipelines.
    \item We propose a two-stage pre-training scheme operating at the patch level that effectively distills geometric representations from noisy, inconsistent part annotations via a multi-positive sample-wise contrastive objective.
    \item We demonstrate that PatchAlign3D achieves fast, single-pass zero-shot part segmentation, significantly outperforming both rendering-based and feed-forward baselines across multiple shape benchmarks.
\end{itemize}
